\part{Finite Groups}
\chapter{Representations of Finite Groups}
\section{Definitions}
\begin{definition}[Representation]\label{dfn:1.1}
	A representation of a finite group $G$ on a finite-dimensional complex vector space $V$ is a group homomorphism $\rho : G \to \operatorname{GL} (V)$ from $G$ to the group $\operatorname{GL} (V)=\Aut_{\Vect{\mathbb{C} } }(V) $. When there is little ambiguity about $\rho $, we simply refer to $V$ as a representation of $G$, and we write $gv$ for $\rho (g)(v)$. We often call the dimension of $V$ the \emph{degree} of $\rho $. We say $\rho $ gives $V$ the structure of a $G$-module.
\end{definition}

If $V,W$ are representations of $G$, then a morphism of representations is a $\mathbb{C} $-linear map $\phi : V \to W$ such that
\[\begin{tikzcd}[row sep = 2.5em, column sep = 2.5em]
	V\arrow[r, "\phi "]\arrow[d, "g"']&W\arrow[d, "g"]\\
	V\arrow[r, "\phi "']&W
\end{tikzcd}\]
commutes for all $g\in G$. That is, $\phi (gv)=g\phi (v)$ for all $g\in G$ and $v\in V$. We call such a map $G$-linear. It follows immediately that $\ker \phi ,\im \phi ,\coker \phi $ are all $G$-modules.

\begin{definition}[Subrepresentation]\label{dfn:1.2}
	A subrepresentation $W$ of a representation $V$ is a subspace $W\subseteq V$ that is $G$-invariant; that is, for all $g\in G$, the restriction $g|_W$ is an element of $\operatorname{GL} (W)$.
\end{definition}

Some common representations are the direct sum $V\oplus W$, given by $g(v,w)=(gv,gw)$, and the tensor product $V\otimes W$, given by $g(v\otimes w)=gv\otimes gw$. This generalizes to tensor powers $V^{\otimes n} $, and we can view the exterior power $\bigwedge^n(V)$ and the symmetric power $\operatorname{Sym}^n(V)$ as subrepresentations of it.

The dual $V^*=\Hom(V, \mathbb{C} ) $ is also a representation. Let $\rho : G \to V$ be a representation, and $\rho ^*: G \to V^*$ be the dual. We want the representations to respect the natural map
\[
	\left\langle \rho ^*(g)(v^*),\rho (g)(v) \right\rangle 
\]
for all $g\in G$, $v\in V$, and $v^*\in V^*$. This definition forces us to define $gv^*$ as the map $(gv^*)(x)=v^*(g^{-1} x)$. The book denotes this as $^t\rho (g^{-1} )$.
\begin{exercise}
	Show that this definition of $\rho ^*$ satisfies the above relation.
\end{exercise}
\begin{proof}
	The fact that the above relation is satisfied is immediate:
	\[
		\left\langle gv^*,gv \right\rangle =(gv^*)(gv)=v^*(g^{-1} gv)=v^*(v)
	.\]
	I feel we must also show $r^*$ is indeed a representation. First, we show $g$ is an automorphism of $V^*$. Surjectivity follows immediately. Let $v^*\in V^*$, and note that since $g$ is an automorphism on $V$, for every $v\in V$ there exists $w\in V$ such that $gw=v$. Therefore, $v^*(gw)=v^*(v)$, and thus the composition
	\[\begin{tikzcd}[row sep = 2.5em, column sep = 2.5em]
		V\arrow[r, "g"]\arrow[rr, bend left, "v^*"]&V\arrow[r, "v^*"]&\mathbb{C} 
	\end{tikzcd}\]
	commutes. Since every $g$ is the inverse of some $g^{-1} \in G$, the given definition is surjective, since $g^{-1} v^*=v^*\circ g\in \Hom(V, \mathbb{C} ) $.

	To show injectivity, we need only show $\ker g=0$ for all $g\in G$. Indeed, if $gv^*=0$, then $v^*(g^{-1} x)=0$ for all $x\in V$. Since $g^{-1} $ is an automorphism of $V$, it follows that $v^*(x)=0$ for all $x\in V$ by surjectivity, and thus $v^*=0$. Therefore, $\rho ^*: G \to \operatorname{GL} V^*$.

	Finally, we show $\rho ^*$ is a group homomorphism. This is immediate; let $g,h\in G$. We have
	\[
		(gh)v^*(x)=v^*\left(\left( h^{-1} g^{-1}   \right) x\right)=(hv^*)\left( g^{-1} x \right) =g(h(v^*))(x)
	.\]
	Therefore, $g\circ h=gh$, and thus $\rho ^*(gh)=\rho ^*(g)\circ \rho ^*(h)$. Therefore, $V^*$ is a representation of $G$ under this definition.
\end{proof}

Recall the following. If $W$ is a finite $k$-vector space, then $W=k^n$ for some $n\in\mathbb{N} $. It follows since Hom is right-adjoint that
\[
	\Hom_{k}(V, W) \cong \Hom_{k}(V, k^n) \cong \Hom_{k}(V, k)^n\cong \left( V^*\otimes_k k \right)^n
.\]
Since tensor is left-adjoint, we have $\left( V^*\otimes _kk \right) ^n\cong V^*\otimes _kk^n\cong V^*\otimes _kW$. Therefore, we would expect $\Hom_{\mathbb{C} }(V, W) $ to be a representation of a group $G$ given the past definitions, provided $V$ and $W$ are representations. Invoking the precise definitions we gave above defines
\[
	(g\phi )(v)=g\left( \phi \left( g^{-1} v \right)  \right) 
\]
for all $\phi \in \Hom(V, W)$. This causes the diagram
\[\begin{tikzcd}[row sep = 2.5em, column sep = 2.5em]
	V\arrow[r, "\phi "]\arrow[d, "g"']&W\arrow[d, "g"]\\
	V\arrow[r, "g\phi "']&W
\end{tikzcd}\]
to commute. That is,
\[
	(g\phi )(gv) = g\left( \phi \left( g^{-1} gv \right)  \right) =g\left( \phi (v) \right) 
,\]
as required. We can in fact view the dual representation as a special case of this, where the $W$ representation is the trivial representation $gw=w$, and thus $(g\phi )(v)=g\phi (g^{-1} v)=\phi (g^{-1} v)$.
\begin{exercise}
	Verify that in general the vector space of $G$-linear maps between two representations $V,W$ is simply the subset $\Hom_{G}(V, W) $ (also denoted $\Hom(V, W) ^G$) of $\Hom(V, W) $ fixed under the action of $G$.
\end{exercise}
\begin{proof}
	I will take ``fixed under the action of $G$'' to mean any $\phi \in \Hom(V, W) $ such that $(g\phi )=\phi $. Note that by our previous definition, for any $v\in V$ and $g\in G$, we ahve
	\[
		(g\phi )(v)=g(\phi (g^{-1} v))
	.\]
	If $\phi $ is $G$-linear, then this equals $\phi (v)$ immediately. Now let $\phi $ be fixed under the action of $G$, and note that this implies
	\[
		g(\phi (g^{-1} v))=\phi (v)\implies \phi (g^{-1} v)=g^{-1} (\phi (v))
	.\]
	Since every group element has a unique inverse, this shows that $\phi (gv)=g\phi (v)$, and $\phi $ is $G$-linear.
\end{proof}

Given the defintions we have here, usual identifications for vector spaces also share the same representations, such as
\[
	\left( V\oplus W \right) \otimes U=\left( V\otimes U \right) \oplus \left( W\otimes U \right) ,
\]
\[
	\extprod^{k}(V\oplus W) =\bigoplus_{a+b=k}\extprod^{a}(V) \otimes \extprod^{b}(V) 
,\]
\[
	\mathbin{\textstyle\bigwedge}^{k}(V^*) =\extprod^{k}(V) ^*
.\]
\begin{exercise}
	Let $\rho : G \to \operatorname{GL} (V)$ be any representation of a finite group $G$ on an $n$-dimensional vector space $V$, and suppose that for any $g\in G$, the determinant $\det \rho (g)=1$. That is, $\rho : G \to \operatorname{SO}_n(\mathbb{C} )$. Show that
	\[
		\extprod^{k}(V) \cong \extprod^{n-k}(V^*) 
	\]
	as $G$-representations.
\end{exercise}
\begin{proof}
	\begin{itemize}
		\item Need to fully realize the precise definition of the $G$-action on exterior powers
		\item Can possibly use the fact that $\extprod^{k}(V) $ is a quotient of $V^{\otimes k} $, and then use
			\[
				\extprod^{n-k}(V) ^*=\Hom(\extprod^{n=k}(V) , \mathbb{C} ) 
			.\]
			Some kind of adjointness condition might be useful.
	\end{itemize}
\end{proof}

Another representation comes from group actions. If $G\to \Aut_{\mathbf{Set} }(X) $ is an action of $G$, then $G$ acts on the free vector space $F^\mathbb{C} (X)$ by
\[
	g\cdot \sum_{i} a_i x_i=\sum_{i} a_i (gx_i)
.\]
The book denotes this by representing each $x\in X$ with a basis vector $e_x$, and giving
\[
	g\cdot \sum a_x e_x=\sum a_x e_{gx} 
.\]
This induces the \emph{regular representation} $R_G$, given by the action of $G$ on itself. This is also sometimes denoted $R$, but this can also mean the space of complex-valued functions $G\to \mathbb{C} $, with action given by $(gf)\left( h \right) =f(g^{-1} h)$.
\begin{exercise}
	Verify that these two descriptions of $R$ agree, by identifying the element $e_x$ with the charateristic function whuch takes the value $1$ on $x$, and $0$ on all other elements of $G$. Show that this representation is then isomorphic to the space of functions on $G$ given by $(g\phi )(h)=\phi (hg)$.
\end{exercise}
\begin{proof}
	This reminds me of coproducts. Consider the map $\phi : R_G \to R$ given by $e_x\mapsto 1_x$, where $1_x$ will denote the function in question. The action of $\phi $ is extended by linearity onto the whole of $R_G$, making it automatically $\mathbb{C} $-linear. Since $G$ is finite, we trivially have an inverse given by $1_x\mapsto e_x$, since these elements are in bijection with $G$. It now suffices to show both $\phi $ and $\phi ^{-1} $ are $G$-linear. By linearity, it suffices to show $g\phi =\phi g$ only on the basis elements, and vice versa. Let $g\in G$. We have
	\[
		\phi (ge_x)=\phi (e_{gx} )=1_{gx}
	.\]
	Note that $1_{gx} (h)=1$ if and only if $h=gx$ if and only if $1_x\left( g^{-1} h \right) =1$. By definiton,
	\[
		1_x\left( g^{-1} \_ \right) =g1_x=g\phi (e_x)
	.\]
	Therefore, $\phi $ is $G$-linear. This same argument can be applied for the reverse direction, showing that these definitions of $R$ are equivalent.

	The space of functions here still denotes the space $G\to \mathbb{C} $, but now under a new action. 
\end{proof}
\section{Complete Reducibility; Schur's Lemma}
We have seen that representations can be constructed with various linear algebraic operations, such as the tensor and direct sum. A classification of representations thus amounts to a classification of representations that are not decomposable into smaller objects, such as direct summands. That is, we look for representations with no proper subrepresentations, often called irreducibles. The following proposition formalizes this notion.

\begin{proposition}\label{prp:1.1}
	If $W$ is a subrepresentation of a representation $V$ of a finite group $G$ over an infinite field $k$, then there is a complementary invariant subspace $W'$ of $V$, such that $V=W\oplus W'$.
\end{proposition}
\begin{proof}
	Since $W\subseteq V$, there exists a complimentary subspace $U$ such that $V=U\oplus W$. This allows us to define the natural projection $\pi_W : V \to W$, and we can then define the average over $G$ of $\pi _W$ as
	\[
		\pi (v)\coloneqq \sum_{g\in G} g(\pi _W(g^{-1} v))
	.\]
	It then suffices to show that $\ker \pi $ is a complimentary subspace of $W$ which is $G$-invariant. For the second statement, it suffices to show $\pi $ is $G$-linear. Indeed, for any $v\in V$, we have
	\[
		\pi (hv)=\sum_{g\in G} g(\pi _W(g^{-1} hv))=h \sum_{g\in G} h^{-1} g\left( \pi _W\left( g^{-1} hv \right)  \right) =h \sum_{g\in G} h^{-1} g\left( \pi _W\left( \left( h^{-1} g \right) ^{-1} v \right)  \right) =h\pi (v)
	.\]
	This last step follows since group multiplication is an automorphism of $G$, by Cayley's theorem. Therefore, $\pi  $ is $G$-linear, and $\ker \pi  $ must be $G$-invariant. Also note that $v\in \ker \pi \cap W$ implies that
	\[
		0=\sum_{g\in G} g(\pi_W(g^{-1} v))=\sum_{g\in G} gg^{-1} v=|G|v,
	\]
	and thus $v=0$ since $\operatorname{char} k=0$. Therefore, $\ker \pi \oplus W=V$.
\end{proof}
\begin{corollary}\label{cor:1.1}
	Any representation is a direct sum of irreducible representations.
\end{corollary}
\begin{proof}
	If a representation $V$ is irreducible, we are done. If not, then there exists a subrepresentation $W$ of $V$, and by proposition \ref{prp:1.1}, there exists a subrepresentation $W'$ of $V$ such that $W\oplus W'=V$.
\end{proof}

We call this property complete reducibility, or semisimplicity. We will see some classes of infinite representations also satisfy this property, with integration wrt an appropriate measure replacing the summation in the averaging part of the proof.

The uniqueness of this decomposition is classified by the following lemma.
\begin{lemma}[Schur]\label{lma:1.1}
	If $V$ and $W$ are irreducible $\mathbb{C} $-representations of $G$ and $\phi : V \to W$ is $G$-linear, then
	\begin{enumerate}
		\item Either $\phi $ is an isomorphism, or $\phi =0$.
		\item If $V=W$, then $\phi =\lambda 1$ for some $\lambda \in\mathbb{C} $, with $1$ the identity map.
	\end{enumerate}
\end{lemma}
\begin{proof}
	The first claim is immediate, as $\ker \phi $ and $\im \phi $ are subrepresentations. For the second, since $\mathbb{C} $ is algebraicly closed, then there exists an eigenvalue $\lambda $ of $\phi $, and thus $\phi -\lambda 1$ has a nonzero kernel. By the first proposition, since $\lambda 1$ is $G$-linear and thus so is $\phi -\lambda 1$, we must have $\phi -\lambda 1=0$, and thus $\phi =\lambda 1$.
\end{proof}
\begin{proposition}\label{prp:1.2}
	For any representation $V$ of a finite group $G$, there exists a decomposition
	\[
		V\cong \bigoplus_{1\leq i\leq k} V_i^{\oplus a_i} 
	,\]
	where the $V_i$s aare distinct irreductible representations. The decomposition is also unique, as are the $V_i$s and their multiplicities $a_i$.
\end{proposition}
\begin{proof}
	Since $V$ is finite, inducting on proposition \ref{prp:1.1} shows the decomposition exists. Suppose $\bigoplus W_i^{\oplus b_i} $ is another representation of $G$ such that $\phi : V \to W$ is an isomorphism. By Schur's lemma, we have that for all $i$ there exists $j$ such that $V_i\cong W_j$ under the restriction $\phi|_{V_i} $. Applying this to the identity $V\to V$  shows the uniqueness criteria.
\end{proof}
Occasionally, the decomposition is written as $\bigoplus a_i V_i$ or $a_1V_i + \cdots + a_n V_n$.

Classifying representations of $G$ thus amounts to the following:
\begin{enumerate}
	\item Classify all irreducible representations of $G$;
	\item Create a formula for the decomposition of $V$ into irreducibles;
	\item Describe all representations that can be obtained from $V$, such as $V\otimes W$, $V^*$, $\Hom(V, W) $, $\extprod^{n}(V) $, and so on.
\end{enumerate}
\section{Examples: Abelian Groups; $\mathfrak{S} _3$}
Representations of abelian groups are easy to work with, largely since the maps induced by a representation $\rho : V \to \operatorname{GL} V$ are in fact $G$-linear homomorphisms. More generally, the map $g(v)$ is $G$-linear if and only if $g\in Z(G)$. It follows by Schur's lemma that if $G$ is abelian, and $V$ is irreducible, then the action of every $g\in G$ on $V$ can be given as a scalar multiple of the identity. It immediately follows that every subspace is invariant simply by common sense, so $V$ must be one dimensional. Therefore, all irreducible representations of abelian groups are 1-dimensional. In other words, all representations $\rho $ must induce maps $\mathbb{C} \to \mathbb{C} $, so
\[
	\rho : G \to \mathbb{C} ^*
\]
for all abelian $G$.

We will now consider the simplest nonabelian group $\mathfrak{S} _3$. We can immediately define two representations on $U$ and $U'$, given by the trivial representation $g(v)=v$, and the \emph{alternating representation} $g(v)=\operatorname{sign} (g)v$. We also have a natural permutation representation on $\mathbb{C} ^3$: if $\left\{ e_1,e_2,e_3 \right\} $ is the standard basis for $\mathbb{C} ^3$, then
\[
	g(v)=g\left( c_1e_1+c_2e_2+c_3e_3 \right) =c_1e_{g(1)} +c_2e_{g (2)} +c_3e_{g (3)} ,
\]
This is \emph{not} irreducible; the subspace generated by $(1,1,1)$ is invariant, and has complimentary subspace $\left\{ (x,y,z)\mid x+y+z=0 \right\} $. This two-dimensional representation is clearly invariant, and we call it the \emph{standard representation of} $\mathfrak{S} _3$.

We now want to classify all representations of $\mathfrak{S} _3$, and we will use a method that will become obsolete for finite groups, but will be useful for Lie groups. If $W$ is a representation of $\mathfrak{S} _3$, then consider the action of the abelian subgroup $\mathfrak{A} _3\subseteq \mathfrak{S} _3$ on $W$. Consider any three-cycle $\tau $, which must generate $\mathfrak{A} _3$. By our previous classification of abelian groups, the representation of $\tau $ is simply an eigenvalue of $\tau $ times the identity matrix; since $\tau ^3=1$, these must be powers of the cube root of unity $\omega =e^{2\pi i/3} $. Thus,
\[
	W=\bigoplus_{i} V_i
,\]
where $V_i=\mathbb{C} v_i$, and $v_i$ ranges over a minimal spanning set of eigenvectors of $\tau $:
\[
	\tau v_i=\omega ^{\alpha _i} v_i
.\]
Now we want to consider how the rest of $\mathfrak{S} _3$ acts on $W$. Consider any transposition $\sigma $ such that $\sigma \tau \sigma =\tau ^2$, and $\sigma ,\tau $ generate $\mathfrak{S} _3$. Note that for an eigenvector $v$ of the action of $\tau $ with eigenvalue $\omega ^i$, we have
\[
	\tau (\sigma (v))=\sigma (\tau ^2(v))=\sigma (\omega ^{2i} v)=\omega ^{2i} \sigma (v)
.\]
Therefore, if $v$ is an eigenvector of $\tau $ with eigenvalue $\omega ^i$, then $\sigma (v)$ is an eigenvector of with eigenvalue $\omega ^{2i} $.

Now choose $v$ an eigenvector of $\tau $ with $\omega^i \neq 1$. Then $\sigma (v)$ is an eigenvector with eigenvalue $\omega ^{2i} \neq \omega ^i$, and thus $\sigma (v)$ is linearly independent of $v$. We can thus take the set spanned by $\left\{ v,\sigma (v) \right\} $, and this will be $\mathfrak{S} _3$-invariant, as any $g\in \mathfrak{S} _3$ is generated by $\tau ,\sigma $. It is shown in the exercises that this space is in fact isomorphic to the standard representation of $\mathfrak{S} _3$.

If the eigenvalue of $v$ under $\tau $ is $1$, then we do not know whether $v$ and $\sigma (v)$ are linearly independent. If not, then $\tau (v)=\tau (\sigma (v))=v$, and the span of $v$ is isomorphic to both the trivial and the alternating representations. If they \emph{are} independent, then $v+\sigma (v)$ and $v-\sigma (v)$ both span one-dimensional representations isomorphic to the trivial  and alternating representations, respectviely.

This shows that there are only three possible irreducible representations of $\mathfrak{S} _3$: standard, alternating, and trivial $V,U',U$. We can thus write
\[
	W\cong U^{\oplus a} \oplus U'^{\oplus b} \oplus V^{\oplus c} 
\]
for any representation $W$ of $\mathfrak{S} _3$. We also have a way to determine $a,b,c$: $c$ is the number of linearly indpendent eigenvectors for $\tau $ with eigenvalue $\omega $, $a+c$ is the multiplicity of $1$ as an eigenvalue of $\sigma $, and $b+c$ is the multiplicity of the eigenvalue $-1$ of $\sigma $.
